%==================================================================
% BAB 1 — PENDAHULUAN
%
% Catatan: teks di bawah adalah CONTOH BERJALAN bertema WebQual 4.0.
% Ganti isinya dengan penelitian Anda, tetapi PERTAHANKAN POLANYA.
%==================================================================

\chapter{PENDAHULUAN}

\section{Latar Belakang}

Program magang merupakan salah satu komponen kurikuler strategis pada Program Studi 
Sistem Informasi Universitas Katolik Indonesia Atma Jaya yang bertujuan untuk 
menjembatani pemahaman teoritis mahasiswa dengan praktik profesional di industri. 
Bagi mahasiswa, pelaksanaan magang bukan sekadar pemenuhan beban SKS, melainkan sarana 
pengembangan kompetensi teknis dan manajerial secara nyata. Namun, efektivitas 
pencapaian tujuan pembelajaran ini sangat bergantung pada tata kelola administratif dan 
operasional yang terstruktur di tingkat program studi.

Pada realitas operasionalnya, alur pelaksanaan magang di Program Studi Sistem Informasi 
Unika Atma Jaya masih dihadapkan pada sejumlah kendala mendasar, khususnya terkait belum 
tersedianya sistem informasi terintegrasi untuk memantau \textit{(monitoring)} setiap tahapan 
alur magang. Proses tata kelola yang saat ini masih dilakukan secara konvensional dan 
manual menyebabkan tingginya kompleksitas koordinasi, ketidakjelasan pelacakan status 
\textit{(status tracking)}, serta kerentanan terhadap ketidaksesuaian prosedur administratif 
antara mahasiswa, pengelola magang, dan dosen pembimbing.

Kondisi serupa juga ditemukan pada penelitian \cite{Rosanti2025} di LP3L Universitas PGRI 
Kanjuruhan Malang, yang mengidentifikasi kesulitan pengelolaan data pendaftaran, ketiadaan 
sistem pemantauan real-time, serta risiko kesalahan input nilai akibat proses yang 
masih manual. Hal yang sejalan juga dilaporkan oleh \cite{Asehan2026}, yang menemukan 
bahwa proses magang berbasis logbook fisik dan dokumen cetak di St. Clare College of 
Caloocan menimbulkan keterlambatan pengumpulan data, ketidakakuratan pencatatan, serta 
minimnya transparansi dalam pemantauan kemajuan magang.

Dampak dari ketiadaan sistem terpusat ini berimplikasi langsung pada aksesibilitas informasi 
dan kejelasan tata laksana magang. Mahasiswa kerap mengalami kesulitan dalam memperoleh 
informasi yang valid dan up-to-date akibat penyebaran sumber informasi yang terfragmentasi. 
Selain itu, ditemukannya ketidakragaman serta ketidakjelasan format dokumen administratif 
magang kerap menimbulkan beban kognitif dan ambiguasi bagi mahasiswa dalam memenuhi persyaratan 
yang ditetapkan. Permasalahan-permasalahan ini tidak hanya menghambat efisiensi operasional 
alur magang, tetapi juga berpotensi menurunkan kualitas pengalaman akademis mahasiswa, 
sebagaimana ditegaskan pula oleh \cite{Asehan2026} yang mencatat bahwa proses manual 
meningkatkan risiko human error, menurunkan akuntabilitas, dan mempersulit penelusuran arsip 
dalam jangka panjang.

Berdasarkan permasalahan tersebut, keberadaan sebuah sistem informasi monitoring alur magang 
yang tersentralisasi dan terstruktur menjadi suatu kebutuhan mendesak. Sistem ini diharapkan 
mampu mengotomatisasi alur kerja, menyediakan transparansi informasi secara real-time, 
mendokumentasikan berkas secara tersentralisasi, serta menyelaraskan standar format dokumen 
administrasi. Pendekatan semacam ini terbukti efektif pada penelitian \cite{Rosanti2025}, 
yang melaporkan tingkat kepuasan pengguna sebesar 88,60 persen setelah penerapan sistem monitoring 
magang berbasis web, maupun pada penelitian \cite{Asehan2026}, yang mencatat peningkatan 
signifikan pada aspek kemudahan pemantauan (skor 4,72) dan kepuasan pengguna secara keseluruhan 
(skor 4,90) setelah penerapan sistem serupa. Dengan demikian, penelitian ini berfokus pada 
pengembangan dan validasi masalah alur proses magang guna menghasilkan solusi berbasis sistem 
informasi yang mampu meningkatkan efektivitas, efisiensi, dan tata kelola pelaksanaan magang 
di Program Studi Sistem Informasi Unika Atma Jaya.

\section{Rumusan Masalah}
% Setiap rumusan masalah diberi \label agar dapat dirujuk dari Bab 3, 4, dan 5.
% Merujuknya cukup dengan \cref{rq:kualitas} -> menghasilkan "RQ1".
% Bila urutannya diubah, seluruh rujukan menyesuaikan sendiri.
Berdasarkan latar belakang di atas, rumusan masalah penelitian ini adalah:

\begin{rumusanmasalah}
    \item Bagaimana merancang sistem informasi yang mampu mengintegrasikan dan memantau 
          (monitoring) setiap tahapan alur magang pada Program Studi Sistem Informasi 
          Unika Atma Jaya secara terpusat?\label{rq:rancang}
    \item Bagaimana membangun sistem yang dapat mengotomatisasi alur kerja administrasi 
          magang guna mengurangi kompleksitas koordinasi antara mahasiswa, pengelola magang, 
          dan dosen pembimbing?\label{rq:otomasi}
    \item Sejauh mana efektivitas dan tingkat penerimaan pengguna terhadap sistem informasi 
          monitoring magang yang dikembangkan?\label{rq:efektivitas}
\end{rumusanmasalah}

\section{Batasan Masalah}
Agar pembahasan terfokus, penelitian ini dibatasi sebagai berikut.

\begin{daftarnomor}
    \item Sistem melibatkan 5 (lima) peran pengguna, yaitu Mahasiswa, Koordinator Magang, 
          Dosen Pembimbing, Pembimbing Eksternal (Tempat Magang), dan Kepala Program Studi.
    \item Ruang lingkup fungsional sistem mencakup proses pengajuan dan validasi pendaftaran 
          magang, pelaksanaan dan bimbingan magang (termasuk pemantauan progres), serta 
          penilaian akhir magang.
    \item Sistem dibangun berbasis aplikasi web (web-based application) yang responsif.
    \item Pengujian sistem dilakukan menggunakan metode \textit{black-box testing}, 
          \textit{User Acceptance Test} (UAT), dan \textit{System Usability Scale} (SUS).
\end{daftarnomor}

\section{Tujuan dan Manfaat}

\subsection{Tujuan Penelitian}
% Tujuan disusun SEJAJAR SATU-SATU dengan rumusan masalah.
% T1 menjawab RQ1, T2 menjawab RQ2, dan seterusnya.
\begin{daftartujuan}
    \item Merancang dan membangun sistem informasi yang dapat memantau setiap tahapan 
          alur magang secara terpusat pada Program Studi Sistem Informasi Unika Atma Jaya.
    \item Mengotomatisasi alur kerja administrasi magang untuk menyederhanakan proses 
          koordinasi antara mahasiswa, pengelola magang, dan dosen pembimbing.
    \item Mengevaluasi tingkat efektivitas dan penerimaan pengguna terhadap sistem yang 
          dikembangkan melalui black-box testing, UAT, dan System Usability Scale.
\end{daftartujuan}

\subsection{Manfaat Penelitian}

\subsubsection{Manfaat Praktis}
Sistem informasi manajemen magang yang dihasilkan dapat digunakan sebagai alat bantu 
untuk memantau setiap tahapan alur magang secara terpusat, mengotomatisasi proses 
administrasi yang sebelumnya dilakukan secara manual, serta menyediakan format dokumen 
yang seragam dan informasi yang transparan bagi mahasiswa. Dengan demikian, beban 
administratif dapat berkurang, koordinasi antarpihak menjadi lebih efisien, dan proses 
pengambilan keputusan terkait pelaksanaan serta penilaian magang dapat dilakukan dengan 
lebih cepat dan akurat.

\subsubsection{Manfaat Akademis}
Penelitian ini diharapkan dapat memberikan kontribusi terhadap pengembangan ilmu 
pengetahuan di bidang sistem informasi, khususnya dalam penerapan konsep perancangan 
sistem berbasis kebutuhan pengguna (user-centered) untuk menyelesaikan permasalahan 
tata kelola magang di lingkungan perguruan tinggi. Penelitian ini juga dapat menjadi 
referensi bagi mahasiswa atau peneliti lain yang tertarik mengembangkan sistem informasi 
monitoring serupa, baik dari sisi metodologi perancangan, pemodelan proses bisnis, 
maupun pengujian sistem menggunakan black-box testing, User Acceptance Test (UAT), 
dan System Usability Scale (SUS). Selain itu, bagi peneliti sendiri, penelitian ini 
menjadi sarana untuk mengaplikasikan dan memperdalam pemahaman teoritis yang diperoleh 
selama perkuliahan ke dalam praktik nyata pengembangan perangkat lunak.

\section{Sistematika Penulisan}
\ifdefstring{\jenisnaskah}{proposal}{%
Proposal Tugas Akhir ini disusun dalam tiga bab dengan rincian berikut.

\noindent \textbf{BAB 1 PENDAHULUAN} memuat latar belakang, rumusan masalah,
batasan masalah, serta tujuan dan manfaat penelitian.

\noindent \textbf{BAB 2 TINJAUAN PUSTAKA} memuat sintesis penelitian terdahulu,
landasan teori yang digunakan, serta justifikasi pemilihan metode.

\noindent \textbf{BAB 3 RANCANGAN PENELITIAN} memuat metodologi penelitian,
populasi dan sampel, instrumen, teknik analisis data, serta rencana dan jadwal
penelitian.
}{%
Penulisan Tugas Akhir ini disusun dalam lima bab dengan rincian berikut.

\noindent \textbf{BAB 1 PENDAHULUAN} memuat latar belakang, rumusan masalah,
batasan masalah, tujuan dan manfaat, serta sistematika penulisan.

\noindent \textbf{BAB 2 TINJAUAN PUSTAKA} memuat sintesis penelitian terdahulu,
landasan teori yang digunakan, serta justifikasi pemilihan metode.

\noindent \textbf{BAB 3 METODE PENELITIAN} memuat tahapan penelitian, populasi
dan sampel, instrumen, teknik analisis data, serta matriks keterkaitan
penelitian.

\noindent \textbf{BAB 4 HASIL DAN PEMBAHASAN} memuat hasil penelitian yang
disajikan per rumusan masalah beserta pembahasannya.

\noindent \textbf{BAB 5 KESIMPULAN DAN SARAN} memuat kesimpulan yang menjawab
setiap rumusan masalah serta saran untuk penelitian selanjutnya.
}
